\section{Introduction}
\IEEEPARstart{I}{ndustrial-agent} deployments must connect each task to its
production context, process signals, and equipment interfaces. Existing work
addresses end-to-end automation control, operator-reviewed control-sequence
generation, modular factory skills, and production scheduling
\cite{xia2025control,huang2026made,koehle2026llmdriven,waseem2026pretrained}.
Other studies target PLC code generation, context-aware control, resource
capability discovery, shop-floor coordination, digital-twin fault handling,
natural-language machine access, and robotic-line control
\cite{liu2026agents4plc,wu2025instructmpc,lim2025llmenabled,zhao2026shopfloor,gill2025llmdt,hofmann2025nlc}.
Across production lines, agent roles, signal maps, protocols, objectives, and
product/run context vary. An integration architecture must accommodate these
differences around shared task and operation services while keeping each
action tied to its active production context.

Consider two production scenes with different process signals and interfaces:
cast-film thickness control and annealing-zone temperature control. An agent
task in either scene needs production context, typed signal access, and a
route from a named operation to the configured device protocol. The framework
must also distinguish a command receipt from a subsequent process
observation. The shared task and operation services can therefore be
configured with scene-specific signal maps, protocols, and limits while
  retaining a common execution path. The evaluation instantiates this path in
  two simulator scenes.

This paper presents AgentWorkShop as a task-to-device integration architecture
and implementation for industrial agents. It connects collaborative tasks
and registered harnesses to production context, manufacturing execution
system (MES) data, analysis and modeling services, data-acquisition (DAQ) and
write-control (DCW) nodes, and protocol-specific drivers. Its access surfaces include command-line, web,
SDK, REST, MCP, and A2A interfaces. Task and chat routes enter the
collaboration runtime, whereas a direct REST tool-invocation route can call
the shared operation dispatcher without creating a task.
Scene blueprints define signal maps, objectives, and guards; channel profiles,
tool policies, and agent bindings scope operations. Agents and callers invoke
named operations through shared runtime checks; the framework does not choose
routes autonomously and retains proposals, command confirmation, and process
observations as distinct records. Fig.~\ref{fig:arch} presents the system
architecture, and Fig.~\ref{fig:runtime} shows the cross-module task flow.
Section~IV details the node-exploration and recipe lifecycles.

The paper makes three contributions:
\begin{enumerate}
\item It specifies a cross-layer task-to-operation contract that assembles
  prompt-time production context, supplies a scoped catalog to a registered
  harness, rechecks each returned call, routes semantic process nodes to
  configured device drivers, and preserves available command and process
  records by route.
\item It documents an implementation architecture that composes collaboration
  tasks, production-context binding, operation checks, semantic I/O, and
  service-owned records. The contribution is this implemented integration
  path; it does not introduce a new agent-reasoning or process-control
  algorithm.
\item It evaluates two simulator paths across cast-film and annealing scenes.
  R20 records six completed, synthetically approved, readback-verified task
  executions with one run-tagged sample per task, all outside target bands.
  Exploratory R13 records protocol and data-path checks while both model
  candidates fail their configured gates. These results characterize the
  tested integration paths rather than process-control efficacy.
\end{enumerate}

Section~II reviews related industrial-agent and integration work.
Section~III presents the design rationale, architecture, and cross-module
execution flow; Section~IV explains module responsibilities and their
integrated operation. Section~V describes the experiments and results, and
Sections~VI and~VII discuss the evidence limits and conclude the paper.
